\paragraph{Elastic Net}

\subparagraph{Ý tưởng}
Elastic Net là phương pháp kết hợp cả hai dạng điều chuẩn $L_1$ (Lasso) và
$L_2$ (Ridge). Mục tiêu là tận dụng ưu điểm của cả hai:
\begin{itemize}
  \item $L_1$: tạo nghiệm thưa (feature selection)
  \item $L_2$: ổn định nghiệm và xử lý tốt khi các biến có tương quan cao
\end{itemize}

\subparagraph{Hạng tử điều chỉnh}
Hạng tử điều chỉnh trong Elastic Net có dạng:
\begin{equation}
  E_W(\boldsymbol{w}) = \alpha \|\boldsymbol{w}\|_1 + \frac{1 - \alpha}{2} \|\boldsymbol{w}\|_2
\end{equation}

\subparagraph{Hàm mục tiêu}
Hàm mục tiêu của Elastic Net là:
\begin{equation}
  E(\boldsymbol{w}) = \frac{1}{2} \|\boldsymbol{X}\boldsymbol{w} -
  \boldsymbol{t}\|^2 + \lambda \left( \alpha \|\boldsymbol{w}\|_1 +
  \frac{1 - \alpha}{2} \|\boldsymbol{w}\|^2 \right)
\end{equation}

\subparagraph{Các trường hợp đặc biệt}
Elastic Net bao gồm hai phương pháp trước như các trường hợp riêng:
\begin{itemize}
  \item $\alpha = 1$: trở thành Lasso Regression
  \item $\alpha = 0$: trở thành Ridge Regression
\end{itemize}

\subparagraph{Nghiệm giải}
Tương tự như Lasso, Elastic Net \textbf{không có nghiệm đóng tường minh} do sự
xuất hiện của chuẩn $L_1$.
Nghiệm thường được tìm bằng các phương pháp tối ưu số:
\begin{itemize}
  \item Coordinate Descent (phổ biến nhất)
  \item Gradient Descent với kỹ thuật proximal
\end{itemize}

\subparagraph{Ưu điểm}
Elastic Net khắc phục được một số hạn chế của Lasso:
\begin{itemize}
  \item Khi các biến đầu vào có tương quan cao:
  \begin{itemize}
    \item Lasso có xu hướng chọn một biến và bỏ các biến còn lại
    \item Elastic Net có thể giữ lại nhóm biến liên quan
  \end{itemize}
  \item Vẫn giữ được tính chất nghiệm thưa (sparse)
  \item Ổn định hơn so với Lasso khi dữ liệu có nhiễu
\end{itemize}

\subparagraph{Diễn giải trực quan}

Elastic Net có thể được hiểu là bài toán với ràng buộc kết hợp:

\begin{equation}
  \min_{\boldsymbol{w}} \|\boldsymbol{X}\boldsymbol{w} - \boldsymbol{t}\|^2
  \quad \text{với} \quad
  \alpha \|\boldsymbol{w}\|_1 + (1-\alpha)\|\boldsymbol{w}\|^2 \leq C
\end{equation}

Elastic Net có thể được hiểu là bài toán tối ưu với ràng buộc:
\begin{equation}
  \min_{\boldsymbol{w}} \|\boldsymbol{X}\boldsymbol{w} - \boldsymbol{t}\|^2
  \quad \text{với} \quad
  \alpha \|\boldsymbol{w}\|_1 + (1 - \alpha)\|\boldsymbol{w}\|^2 \leq C
\end{equation}

Miền ràng buộc trong trường hợp này là sự kết hợp giữa hình thoi ($L_1$) và
hình cầu ($L_2$).

Do đó:
\begin{itemize}
  \item Giống Lasso: nghiệm có xu hướng thưa (nhiều hệ số bằng 0)
  \item Giống Ridge: nghiệm ổn định hơn và ít nhạy với nhiễu
\end{itemize}

Elastic Net tạo ra sự cân bằng giữa việc lựa chọn đặc trưng và duy trì tính ổn
định của mô hình, đặc biệt hiệu quả khi các biến đầu vào có tương quan cao.
